\section{Limitations and Ethical Considerations}
\label{sec:limitations}

Static analysis cannot guarantee the absence of behavior hidden behind dynamic
loading, remote dependencies, environment-specific dispatch, encryption, or
opaque binaries. The graph currently provides AST-level flow only for Python;
other languages retain lexical evidence, and the labeled corpora do not contain
enough malicious Python flow to measure path recall. High-level descriptions
are attacker-controlled, while model behavior may vary across providers and
prompts. Consequently, \textsc{Pass} is not a proof of safety.

The source-disjoint evaluation is balanced and therefore does not estimate
real-world prevalence. Thirteen Ours samples and ten direct-model samples remain
unresolved after bounded retries; primary comparisons use the 977-package
intersection rather than imputing outcomes. The current-community scan retains
only primary Skill documents and has no truth labels. Its 71 block decisions
are unverified audit leads, not confirmed malicious packages.

Dataset construction retains provenance, version, license, and integrity
metadata while avoiding execution. Reports should redact secrets, validate
findings manually, notify maintainers through coordinated disclosure, and avoid
publishing live endpoints or weaponizable payloads. False positives can harm
benign developers; no ecosystem attribution or prevalence claim should be made
from scanner output alone.
